\chapter{Virasoro Representations, Minimal Models, and BPZ Equations}
\label{ch:volume-iii-virasoro-representations-bpz}

\section{Highest-Weight Modules}

A chiral Virasoro highest-weight state \(\ket{h}\) is defined by
\[
  L_0\ket{h}=h\ket{h},
  \qquad
  L_n\ket{h}=0\quad(n>0).
\]
Descendants are generated by negative modes:
\[
  L_{-n_1}L_{-n_2}\cdots L_{-n_k}\ket{h},
  \qquad
  n_i>0.
\]
The level is
\[
  N=n_1+\cdots+n_k,
\]
and the \(L_0\) eigenvalue of a level-\(N\) descendant is \(h+N\).

Unitarity gives an inner product with
\[
  L_n^\dagger=L_{-n}.
\]
At level one,
\[
  \|L_{-1}\ket{h}\|^2=2h\braket{h}{h},
\]
so \(h\ge0\).  At level two, in the basis
\[
  L_{-2}\ket{h},
  \qquad
  L_{-1}^2\ket{h},
\]
the Gram matrix is
\begin{equation}
  M^{(2)}
  =
  \begin{pmatrix}
    4h+\frac c2 & 6h\\
    6h & 4h(2h+1)
  \end{pmatrix}.
  \label{eq:virasoro-level-two-gram}
\end{equation}
Positivity of all Gram matrices gives the full Virasoro unitarity
constraints.

\section{Null States}

A null state is a descendant with zero norm that is orthogonal to the entire
module.  In the physical Hilbert space, it is quotiented out.  In local
operator language, a null state gives a differential equation for correlation
functions.

At level two, a null state can take the form
\begin{equation}
  \left(
    L_{-2}
    -
    \frac{3}{2(2h+1)}L_{-1}^2
  \right)\ket{h}=0,
  \label{eq:level-two-null-vector}
\end{equation}
with \(h\) and \(c\) related by the vanishing of the determinant of
\eqref{eq:virasoro-level-two-gram}.  The precise relation is encoded by the
Kac parametrization below.

\section{Unitary Minimal Models}

The unitary Virasoro representations occur in two branches.  One branch has
\[
  c\ge1,
  \qquad
  h\ge0.
\]
The discrete branch is the minimal-model series
\begin{equation}
  c_m=1-\frac{6}{m(m+1)},
  \qquad
  m=3,4,5,\ldots .
  \label{eq:unitary-minimal-central-charge}
\end{equation}
The allowed highest weights are
\begin{equation}
  h_{r,s}^{(m)}
  =
  \frac{\bigl((m+1)r-ms\bigr)^2-1}{4m(m+1)},
  \qquad
  1\le r\le m-1,
  \quad
  1\le s\le m,
  \label{eq:kac-weights-unitary-minimal}
\end{equation}
with the identification
\[
  (r,s)\sim(m-r,m+1-s).
\]
The full local CFT is obtained by combining holomorphic and antiholomorphic
representations subject to locality and modular consistency.

\begin{figure}[H]
  \centering
  \resizebox{0.85\textwidth}{!}{%
  \begin{tikzpicture}[x=1cm,y=1cm,every node/.style={font=\scriptsize}]
    \tikzset{axis/.style={-{Latex[length=2mm]},line width=0.8pt}}
    \draw[axis] (-0.2,0)--(6.4,0) node[right] {\(c\)};
    \draw[axis] (0,-0.3)--(0,3.1) node[above] {\(h\)};
    \draw[blue!65!black,line width=1.0pt] (3.15,0)--(6.15,0);
    \node[blue!65!black] at (5.2,0.35) {\(c\ge1,\ h\ge0\)};

    \foreach \x/\lab in {0.85/{m=3},1.45/{4},1.90/{5},2.25/{6},2.55/{\cdots}} {
      \draw[magenta,line width=0.9pt] (\x,0.0)--(\x,2.15);
      \node[magenta,below] at (\x,0) {\(\lab\)};
    }
    \node[magenta] at (1.75,2.55) {minimal models};
    \draw[dashed] (3.0,0)--(3.0,2.65);
    \node at (3.0,-0.28) {\(1\)};
  \end{tikzpicture}
  }
  \caption{The unitary Virasoro highest-weight representations consist of the
  continuous branch \(c\ge1\) and the discrete minimal-model series
  \(c_m<1\).}
  \label{fig:virasoro-unitarity-branches}
\end{figure}

\section{BPZ Differential Equations}

Let \(\phi_h(z)\) be a chiral primary with a level-two null state
\eqref{eq:level-two-null-vector}.  In a correlation function with other
chiral primaries \(\mathcal O_i(z_i)\), the insertion of the null state
vanishes:
\[
  \left\langle
    \left(
      L_{-2}
      -
      \frac{3}{2(2h+1)}L_{-1}^2
    \right)\phi_h(z)
    \prod_i\mathcal O_i(z_i)
  \right\rangle
  =
  0.
\]
The \(L_{-1}\) insertion is \(\partial_z\phi_h\).  The \(L_{-2}\) insertion is
computed by surrounding \(\phi_h(z)\) with a stress-tensor contour and moving
that contour to the other insertions.  This gives the BPZ equation
\begin{equation}
  \left[
    \frac{3}{2(2h+1)}\partial_z^2
    -
    \sum_i
    \left(
      \frac{h_i}{(z_i-z)^2}
      -
      \frac{1}{z_i-z}\partial_{z_i}
    \right)
  \right]
  \left\langle
    \phi_h(z)\prod_i\mathcal O_i(z_i)
  \right\rangle
  =
  0.
  \label{eq:bpz-level-two-equation}
\end{equation}

\begin{figure}[H]
  \centering
  \resizebox{0.92\textwidth}{!}{%
  \begin{tikzpicture}[x=1cm,y=1cm,every node/.style={font=\scriptsize}]
    \tikzset{
      insertion/.style={circle,fill=violet!80!black,inner sep=1.7pt},
      contour/.style={green!45!black,line width=0.85pt},
      arrow/.style={-{Latex[length=2mm]},line width=0.8pt}
    }
    \begin{scope}
      \node[font=\small] at (0,1.95) {stress-tensor contour};
      \draw[thick] (-1.95,-1.25) rectangle (1.95,1.25);
      \node[insertion] at (-0.45,0.0) {};
      \node[violet!80!black] at (-0.55,0.35) {\(\phi_h(z)\)};
      \foreach \x/\y/\lab in {0.85/0.55/{z_1},0.75/-0.55/{z_2},-1.30/0.55/{z_3}} {
        \node[insertion] at (\x,\y) {};
        \node[violet!80!black,above right] at (\x,\y) {\(\mathcal O(\lab)\)};
      }
      \draw[contour] (-0.45,0) circle (0.36);
    \end{scope}

    \draw[arrow] (2.45,0)--(3.65,0);
    \node[align=center] at (3.05,0.52) {move contour\\to other insertions};

    \begin{scope}[shift={(6.25,0)}]
      \draw[thick] (-1.95,-1.25) rectangle (1.95,1.25);
      \foreach \x/\y/\lab in {0.85/0.55/{z_1},0.75/-0.55/{z_2},-1.30/0.55/{z_3}} {
        \node[insertion] at (\x,\y) {};
        \draw[contour] (\x,\y) circle (0.32);
      }
      \node[blue!65!black] at (0,-1.55) {residues become differential operators};
    \end{scope}
  \end{tikzpicture}
  }
  \caption{A null state produces a differential equation because the
  stress-tensor contour can be moved from the degenerate insertion to the other
  insertions.}
  \label{fig:bpz-contour}
\end{figure}

\section{The Ising Minimal Model}

The first unitary minimal model has \(m=3\) and central charge
\[
  c=\frac12.
\]
Its primary weights are
\[
  h_{1,1}=0,
  \qquad
  h_{2,1}=\frac12,
  \qquad
  h_{1,2}=h_{2,2}=\frac1{16}
\]
up to the Kac-table identification.  The diagonal local theory has three
primary fields:
\[
  \mathbf1,
  \qquad
  \epsilon\quad (h,\bar h)=\left(\frac12,\frac12\right),
  \qquad
  \sigma\quad (h,\bar h)=\left(\frac1{16},\frac1{16}\right).
\]
The spin field obeys the null-vector relation
\begin{equation}
  \left(L_{-2}-\frac43L_{-1}^2\right)\ket{\sigma}=0,
  \label{eq:ising-sigma-null}
\end{equation}
and the same relation in the antiholomorphic sector.

For the normalized four-point function of spin fields,
\[
  \left\langle
    \sigma(z_1,\bar z_1)\sigma(z_2,\bar z_2)
    \sigma(z_3,\bar z_3)\sigma(z_4,\bar z_4)
  \right\rangle
  =
  |z_{12}|^{-2\Delta_\sigma}|z_{34}|^{-2\Delta_\sigma}f(z,\bar z),
\]
where
\[
  \Delta_\sigma=\frac18,
\]
the BPZ equation reduces, in the standard conformal frame, to
\begin{equation}
  \left[
    \partial_z^2
    +
    \frac{2-5z}{4z(1-z)}\partial_z
    -
    \frac{3}{64(1-z)^2}
  \right]f(z,\bar z)=0,
  \label{eq:ising-sigma-bpz}
\end{equation}
with an analogous equation in \(\bar z\).  Two holomorphic solutions are
\[
  f_\pm(z)=(1-z)^{-1/8}\sqrt{1\pm\sqrt z}.
\]
Single-valuedness and normalization determine
\begin{equation}
  f(z,\bar z)
  =
  \frac{|1+\sqrt z|+|1-\sqrt z|}
       {2|1-z|^{1/4}}.
  \label{eq:ising-spin-four-point}
\end{equation}
Comparing with the Virasoro block expansion gives
\[
  C_{\sigma\sigma\epsilon}=\frac12
\]
in the normalization where the two-point functions of primaries are one.  The
\(\mathbb Z_2\) selection rule gives
\[
  C_{\epsilon\epsilon\sigma}=0.
\]
